% 6. Limitations
\section{Limitations}

\begin{enumerate}
\item \textbf{Structural, not functional.} All metrics derive from the
static wiring diagram. Node removal is a computational abstraction of
silencing; no physiological manipulation was performed, and no claim of
functional necessity is made.
\item \textbf{Single reconstruction.} FAFB v783 is one female brain;
reconstruction- and specimen-dependence is unquantified here.
\item \textbf{Metric dependence.} Global directed efficiency is one
functional; the reachability-drop diagnostic correlates only moderately
($\rho=0.39$) and both are reported.
\item \textbf{Neurotransmitter annotations.} NT identities are
model predictions \cite{schlegel2024}, not measured signs; GLUT is
sign-ambiguous and was analyzed separately.
\item \textbf{Degree matching.} Greedy 1:1 matching within $\pm10\%$
cannot exclude residual within-window imbalance; the exact
degree-preserving null ensemble addresses the dominant confound.
\item \textbf{Connection-table sensitivity not completed.} The full
no-threshold sensitivity analysis was RAM-infeasible in the 8-GB
environment ($\sim$50M collapsed pairs; out-of-memory), and is deferred
to a high-RAM notebook; the Buhmann table is methodologically
non-comparable and was excluded. This analysis is \emph{not} claimed as
performed.
\item \textbf{Null-resolution and estimator matching.} With $n=100$
nulls, empirical tail resolution is ${\sim}0.01$; the null comparison
uses the $k{=}4$ panel for observed and nulls alike, while the primary
matched-pair statistic reads $\delta=0.111$ from the $k{=}8$ screening
CIS. Both pre-registered statistics agree on the verdict.
\end{enumerate}
